\documentclass[11pt]{article}
\usepackage{amsmath,amssymb,amsthm,hyperref}
\newtheorem{theorem}{Theorem}
\newtheorem{lemma}[theorem]{Lemma}
\newtheorem{corollary}[theorem]{Corollary}
\newcommand{\Q}{\mathbb Q}
\newcommand{\Z}{\mathbb Z}
\newcommand{\C}{\mathbb C}
\newcommand{\CH}{\operatorname{CH}}
\newcommand{\rk}{\operatorname{rank}}
\title{Algebraizing a restored GST fold: exact torsor criterion and actual cycle return}
\author{Hand derivation before Lean; author self-review}
\date{3 November 2026}
\begin{document}\maketitle

\section{The exact subconstruction and original target}
Base checkpoint: \texttt{e699dbec256fc23d4c5615449bf303df45d0076d}.
The original input remains arbitrary smooth projective complex $X$, integer
$1\le p\le\dim X$, and prescribed rational $(p,p)$ class $\alpha$.
The desired final output is $Z_\alpha\in\CH^p(X)_\Q$ with canonical
Betti class $\alpha$. This paper executes the intervening projective
algebraization subconstruction for the actual elliptic-torus folds over $X$.
It does not identify a rational homotopy-fiber model in higher degree with
one of these principal torus bundles.

Work first on connected $X$; all statements extend componentwise. Fix
$E=\C^*/q^\Z$, $q=e^{-2\pi}$, and $A=E^r$. For actual holomorphic (hence
algebraic) lines $L_i$ on $X$ let
\[
 P=\prod_X (L_i^\times/q^\Z),\quad
 \theta=(\theta(L_i))\in H^1(X,\mathcal A),\quad
 C:\Q^r\to H^2(X,\Q),\ e_i\mapsto c_1(L_i).
\]
Here $\mathcal A$ is the analytic sheaf of holomorphic maps to $A$;
$H^1$ classifies principal $A$-torsors with their specified action over $X$.
The prior curvature proof establishes whole canonical Hodge purity iff
$C=0$. It also retains the distinct gluing class $\theta$.
No Hodge-conjecture premise is used below.

\section{The projectivity test, including both directions}
\begin{theorem}[Exact projective algebraization criterion]\label{thm:proj}
For every holomorphic principal $A$-torsor $\pi:P\to X$ with $X$ smooth
projective and $A$ a fixed projective complex torus, the following are equivalent:
\begin{enumerate}
\item $P$ is a projective complex manifold;
\item $\theta$ is torsion: $[N]\theta=0$ for some integer $N>0$;
\item an actual finite unramified cover $f:X'\to X$ trivializes $P$, and
$P$ has the finite quotient presentation $(X'\times A)/A[N]$ constructed below.
\end{enumerate}
When these hold, $P$, $\pi$, and the principal action have algebraic
realizations over the original $X$.
\end{theorem}
\begin{proof}
Suppose $P$ is projective and take an ample line $M$ on it. Trivialize $P$
on small analytic opens $U_i$ by choices of local origins. On each fiber,
$M$ has an ample Neron--Severi class $\lambda$, constant on connected $X$.
Choose an actual ample line $M_0$ on $A$ in that class. The fiber restrictions
of $M$ define holomorphic maps $a_i:U_i\to\operatorname{Pic}^0(A)=A^\vee$
after tensoring by $M_0^{-1}$. Changing the origin by transition $t_{ij}$
changes those maps by $\lambda_{M_0}(t_{ij})$, up to the fixed harmless
choice of sign, where
\[
 \lambda_{M_0}:A\to A^\vee,\quad
 t\mapsto T_t^*M_0\otimes M_0^{-1}.
\]
Thus $\lambda_{M_0,*}\theta=0$ in $H^1(X,\mathcal A^\vee)$.
Positivity of $M_0$ makes $\lambda_{M_0}$ an isogeny: its positive Hermitian
Riemann form is nonsingular, so its induced map on the real universal covers
is invertible and its lattice cokernel is finite. Choose $N>0$ annihilating
that finite cokernel; the inverse on universal covers then gives an isogeny
$\mu:A^\vee\to A$ with $\mu\lambda_{M_0}=[N]$. Consequently $[N]\theta=0$.
This argument uses an actual ample line on the total space, not purity alone.

Conversely, the exact analytic sheaf sequence
\[
 0\longrightarrow A[N]\longrightarrow\mathcal A
 \xrightarrow{[N]}\mathcal A\longrightarrow0
\]
is locally surjective because $[N]$ is a covering map. If $[N]\theta=0$,
its long exact sequence gives an actual $A[N]$-valued cocycle lifting
$\theta$. Its principal cover $X'\to X$ is finite unramified of degree
$d=|A[N]|=N^{2r}$ (it need not be connected). The associated bundle is
\[
 P\cong(X'\times A)/A[N],
 \qquad g\cdot(x',a)=(g x',a-g).
\]
Local cocycles verify this equality with the original torsor, not just
with its topological class. Pullback of a positive ample line from $X$
is positive on compact $X'$, so Kodaira embedding makes every component
projective. The covering is algebraic by projective GAGA. The finite group
actions are algebraic for the same reason. A finite group quotient of a
projective variety is projective: tensor an ample line over the group orbit,
then take a sufficiently high power for descent and invariant sections.
The action here is free, so the quotient is smooth. The quotient map to $X$
and the translation action are algebraic. This constructs (3) and (1).
The implication (3) to (2) also follows from the displayed finite cocycle.
\end{proof}
The analytic inputs used are the line-family map to $\operatorname{Pic}^0$
and the Kodaira/projective-GAGA facts at their usual precise strengths.
No statement that all K\"ahler manifolds are projective is used.

\section{Actual multisection cycle and normalized canonical return}
\begin{theorem}[Source constructed by the finite cover]\label{thm:cycle}
Under the equivalent conditions above, set $s:X'\to P$ to be the algebraic
map induced by $x'\mapsto(x',0)$. With $d=N^{2r}$ the genuine rational
Chow class
\[
 \Sigma=\frac1d s_*[X']\in\CH^r(P)_\Q
\]
satisfies $\pi_*\Sigma=[X]$. Therefore
\[
 R_\Sigma(B)=\pi_*(\Sigma\cdot B):\CH^p(P)_\Q\to\CH^p(X)_\Q,
 \qquad R_\Sigma\pi^*=\mathrm{id}.                 \tag{1}
\]
These statements hold in actual Chow, with the canonical Betti compatibility.
\end{theorem}
\begin{proof}
The map $s$ is proper and finite onto its image, since its composite with
$\pi$ is finite and its fibers are finite. All maps have already been
algebraized. Hence
$\pi_*s_*[X']=f_*[X']=d[X]$, including disconnected covers and multiplicities.
The projection formula gives
$\pi_*(\Sigma\cdot\pi^*Z)=Z\cdot\pi_*\Sigma=Z$.
Smoothness defines the Chow intersection; proper pushforward lowers
codimension by $r$, so (1) has the stated degree $p$. Compatibility of the
canonical cycle map with these operations gives the Betti identities.
\end{proof}
In particular, for the actual prescribed $\alpha$,
\[
 \alpha\text{ is a cycle class on }X
 \iff \pi^*\alpha\text{ is a cycle class on }P.       \tag{2}
\]
The forward map is actual pullback; in the reverse direction an actual
$B$ with $\mathrm{cl}(B)=\pi^*\alpha$ yields the explicit
$Z_\alpha=R_\Sigma(B)$, with class $\alpha$. The construction of $\Sigma$
is unconditional once the torsor is torsion; it does not construct $B$
for an arbitrary higher-primitive $\alpha$. Equation (2) prevents silently
using the provisional Hodge assumption at the projective endpoint.

The finite cover also computes the entire rational Hodge cup algebra.
Translations on $A$ act trivially on cohomology, being isotopic to identity.
Finite-cover transfer and Kunneth give
\[
 H^*(P,\Q)\cong
 (H^*(X',\Q)\otimes H^*(A,\Q))^{A[N]}
 \cong H^*(X,\Q)\otimes H^*(A,\Q).                 \tag{3}
\]
These are Hodge-algebra isomorphisms because all maps are algebraic and
finite invariants are exact over $\Q$. The splitting uses the finite-cover
presentation and is not claimed canonical under all choices. Every degree
is retained; no first-Betti surrogate or claim about the full Chow ring is
being made.

\section{Largest projective dimensional quotient}
Return to $A=E^r$ and the list of lines. Define three actual subgroups of
$\Z^r$:
\[
\begin{split}
 R_0&=\{m:\textstyle\sum_i m_i\theta(L_i)=0\},\\
 R_{\rm alg}&=\{m:\textstyle\sum_i m_i\theta(L_i)\text{ is torsion}\},\\
 R_C&=\{m:\textstyle\sum_i m_i c_1(L_i)=0\text{ in }H^2(X,\Q)\}.
\end{split}
\begin{theorem}[Exact projective quotient and dimension gap]\label{thm:max}
In the category of integer-matrix principal $E$-quotients over $X$,
\[
 R_{\rm alg}=\operatorname{sat}(R_0),\qquad
 R_0\subset R_{\rm alg}\subset R_C.
\]
In particular $R_{\rm alg}$ is saturated. Choose an integral basis as the
rows of $M_{\rm alg}$. The corresponding connected-fiber quotient has a
projective total space, dimension $n+s$ where $s=\rk R_{\rm alg}$, and
every projective integer-matrix quotient factors uniquely through it.
The sharp number of dimensions that must collapse for projectivity is
$r-s$. The extra collapse beyond Hodge restoration is
\[
 \Delta_{\rm ext}=\rk R_C-\rk R_{\rm alg}\geq0.       \tag{4}
\]
\end{theorem}
\begin{proof}
A relation is torsion iff a positive multiple is zero; this proves the
saturation equality and closure under sums, including different torsion
orders. A torsion torsor has torsion integral topological Chern class, so
its rational curvature vanishes; hence $R_{\rm alg}\subset R_C$.
A saturated subgroup of a finite free lattice has free quotient; extend
its basis to a unimodular basis of $\Z^r$. Thus the basis-row map is
surjective on lattices and its map $E^r\to E^s$ has connected kernel.
Its finitely many torsor coordinates are torsion. A common multiple
annihilates all of them, so Theorem~\ref{thm:proj} constructs its projective
endpoint and Theorem~\ref{thm:cycle} its multisection.
Every projective quotient has torsion coordinate classes, so all its rows
lie in $R_{\rm alg}$. Unique integral expansion in the basis gives the
factorization. If the original row map is surjective, so is this factor
map, since a right inverse for $M_{\rm alg}$ is integral. Its row rank
cannot exceed $s$, proving optimality. The earlier Hodge maximum is
$n+\rk R_C=n+r-\rk C$; subtraction proves (4).
For $r=0$ the empty quotient is just $X$, with $N=1$ and $\Sigma=1$.
\end{proof}
The optimum concerns precisely these integer-matrix quotients. It is not
an optimality claim over all GST realizations. Basis choices are not
canonical. Computing the analytic torsor classes is not replaced by a
claimed decision procedure from curvature alone.

\section{A geometric discriminator, rather than another rank extrapolation}
If $h^{1,0}(X)>0$, the earlier countability argument gives a line
$L\in\operatorname{Pic}^0(X)$ for which $\theta(L)$ is not torsion.
Indeed the image of $H^1(X,\Z)\to\operatorname{Pic}(X)$ under $q^{(-)}$
is countable, as is the union of its preimages under all nonzero power
maps on the positive-dimensional compact torus $\operatorname{Pic}^0(X)$;
each such power map has finite fibers. Choose $L$ outside that union.
Then $P(L)$ has $C=0$, its entire canonical Hodge structure and the actual
K\"ahler metric from the prior construction. Yet Theorem~\ref{thm:proj}
proves it has no projective realization of the given complex manifold.
Here $R_C=\Z$, $R_{\rm alg}=0$: the maximal Hodge quotient keeps the handle,
and the maximal projective quotient collapses it, returning to $X$.
This refutes the proposed implication ``curvature zero automatically
algebraizes every restored fold.'' It does not refute the Hodge conjecture.

For the earlier independent central worlds $W_K=X\times Z_K$, zero total
curvature gives the actual product $X\times E^{m+r}$; its torsor is zero.
Thus that experiment lies on the positive branch, with the explicit
multisection and original-X return above. Its curvature theorem remains
established. The new test concerns the broader base-coupled construction.

\section{Executed routes and remaining source scope}
Three distinct checks were executed: (i) construct a projective model from
a torsion gluing class, with finite quotient and actual multisection;
(ii) derive the necessary torsion condition from an actual ample line,
then apply it to the already constructed zero-curvature non-torsion world;
(iii) test the proposed alpha return using the actual multisection, obtaining
the exact two-way cycle-image relation (2), rather than an assumed source.
The projective realization subconstruction and its optimal quotient are
completed at hand-proof level, self-reviewed. No worldwide novelty claim
is made for the classical torsor/projectivity ingredients.
For higher $p$ the prescribed rational homotopy-fiber torsor is not a
principal elliptic torsor; these theorems do not silently algebraize it.
The universal $B$ in (2), or a different independently algebraic higher
source representing $\alpha$, has not been constructed by this attempt.
\end{document}
